\section{Esempi}

\subsection{Esempio di tabella}


\begin{table}
    \centering
    \begin{tabular}{| m{1.8cm} ||m{1.8cm}|m{1.8cm}|m{1.8cm}|m{1.8cm}|m{1.8cm}|m{1.8cm}|}
    \hline
         &  fine grained & facetted & featureless & smooth-facetted & facetted-oriented & large facetted\\
    \hline\hline
        fine grained      & 15 & 0 & 0 & 1 & 0 & 0\\
        \hline
        facetted          & 0  & 12 & 0 & 0 & 0 & 0 \\[3ex]
        \hline
        featureless       & 0 & 0 & 1 & 1 & 0  & 0\\[3ex]
        \hline
        smooth-facetted   & 0 & 0 & 0 & 5 & 0 & 0 \\
        \hline
        facetted-oriented & 0 & 0 & 0 & 0 & 13 & 0\\
        \hline
        large facetted    & 0 & 0 & 0 & 0 & 0  & 2\\
        \hline
         
    \end{tabular}
    \caption{Esempio di una tabella normale.}
    \label{tab:riferimento a tabella}
\end{table}

\begin{table}[hb]
    \begin{tabularx}{\textwidth} {  | >{\raggedright\arraybackslash}X 
                                    || >{\centering\arraybackslash}X 
                                    | >{\centering\arraybackslash}X | }
        \hline
        & Community Hub & Business Hub \\
        \hline\hline
        Collaborazione & & \\
        \hline
        Navigazione download di componenti, estensioni e workflows & \ding{51} & \ding{51} \\
        \hline
        Spazio privato & \ding{51} & \ding{51} \\
        \hline
        Collaborare ai workflow & Solo spazi pubblici o pacchetto Team & \ding{51} \\
        \hline
        Accesso il lettura per gli utenti non registrati & \ding{51} & Solo pacchetto Enterprise \\
        \hline
        Tracciamento versioni & \ding{55} & \ding{51} \\
        \hline
        Collezioni & \ding{55} & Solo pacchetto Enterprise \\
        \hline\hline
        Automazione &  & \\
        \hline
        Esecuzione nell'Hub & \ding{55} & \ding{51}\\
        \hline
        Esecuzione programmata & \ding{55} & \ding{51}\\
        \hline
        Esecuzione condizionale & \ding{55} & \ding{51}\\
        \hline
        Vcores e ambiente dedicato e scalabile & \ding{55} & \ding{51}\\
        \hline \hline
        Distribuzione & & \\
        \hline
        Data Apps & \ding{55} & \ding{51}\\
        \hline
        REST API & \ding{55} & \ding{51}\\
        \hline
        Knime Edge & \ding{55} & \ding{51}\\
        \hline

        
    \end{tabularx}
    \caption{Tabella con tabularx.}
    \label{tab:esempio tabularx}
\end{table}


Esempio di tabella normale a \ref{tab:riferimento a tabella}.

Esempio di tabella con tabularx a \ref{tab:esempio tabularx}.

\subsection{Esempio di codice}

\begin{lstlisting}[style=py, caption=Script Python.]
#for loop that traverses numbers from 1 to 100
for i in range(1,101):
  #check if number is divisible by both 3 and 5
  if(i%3==0 and i%5==0):
    print("FizzBuzz")
  #check if number is divisible by 3
  elif(i%3 == 0):
    print("Fizz")
  #check if number is divisible by 5
  elif(i%5 == 0):
    print("Buzz")
  #if not divisible by either of them print the i
  else:
    print(i)
\end{lstlisting}
\label{script}

Volendo si possono anche mettere note a piè di pagina\footnote{Questa è una nota a piè pagina} oppure inserire link \url{https://www.youtube.com/watch?v=G1IbRujko-A}.